\section{System 1: Structure of the Substrate}
\label{sec:system1_projection}
With the substrate fixed as $E_8$, the remaining postulates of \Cref{tab:postulates} fix its internal structure: sectors, channels, signature, boundary, and clock. Together they force every integer used in System 2 (Lean: \texttt{Derived.parameters\_forced}).

\subsection{Sectors}
\label{sec:symmetric_decomposition}
\textbf{Postulate B.} The substrate contains, with finite index, two orthogonal rank-four sectors of the same type, one for addresses and one for internal state. Each is generated by norm-two vectors, is irreducible (not itself an orthogonal sum), and is bit-addressed ($2\Lambda_s^*\subseteq\Lambda_s$, Postulate A applied to a sector).

A positive-definite lattice generated by norm-two vectors is a simply-laced root lattice, so the rank-four candidates are $A_4$, $D_4$, $A_3A_1$, $A_2A_2$, $A_2A_1A_1$, and $A_1^4$. Only $D_4$ is irreducible with involutive glue: $A_4$ has glue group $\mathbb Z/5$ and $A_1^4$ is reducible (Lean: \texttt{sector\_selection}). $E_8$ contains two orthogonal $D_4$ sectors, given by an explicit basis whose Gram matrix is $D_4\oplus D_4$ (\texttt{e8\_two\_d4\_sectors}); their sum has index four, and its four glue cosets reconstruct $E_8$. The address sector has rank $D=4$.

\subsection{Channels}
\label{sec:chiralcapacity}
A sector's glue group $D_4^*/D_4\cong(\mathbb Z/2)^2$ has three nonzero classes. We define a \emph{channel} as a nonzero glue class, and its capacity as its number of minimal (least-cost) vectors. Each of the three classes has exactly eight (Lean: \texttt{minimal\_class\_card}); these are the classes usually labeled $8_v$, $8_s$, $8_c$ and permuted by triality\cite{slansky_group_1981}.

Directed flow needs a source and a sink, so at least two channels (Causality); Economy allows no more. Hence each sector carries
\begin{equation}
\nu=2\times8=16
\end{equation}
channel coordinates, whichever two classes are used, and the node carries $N=2\nu=32$ (\texttt{channel\_total}). The third class carries no flow. Pure lattice symmetry would treat all three classes alike; the two-channel rule is where Causality breaks triality.

\subsection{Signature and Boundary}
\label{sec:metric_signature}
For a four-dimensional orthogonal Clifford basis with $q$ negative squares, the volume element obeys $\omega^2=(-1)^q$ (Lean: \texttt{signature\_volume\_square}).

\textbf{Postulate G.} The volume element of the address sector supplies its complex structure, $\omega^2=-1$; with the overall sign convention, time directions are the minority.

Then $q=1$: one time direction and three spatial ones (\texttt{one\_time}). The internal configuration cost remains positive definite; the Lorentzian signature belongs to the effective coordinate sector.

\label{sec:topological_boundary}
\textbf{Postulate H.} A persistent defect is bounded by a closed connected surface of no greater topological complexity than any other, that is, of maximal Euler characteristic.

Among closed surfaces $\chi\le2$, with equality only for the sphere (Lean: \texttt{euler\_eq\_two\_iff}, \texttt{boundary\_sphere}). So the boundary is $S^2$ and
\begin{equation}
\chi=2.
\end{equation}
A defect's configuration is a spatial position and a boundary point, of dimension
\begin{equation}
\sigma=(D-q)+2=3+2=5.
\end{equation}

\subsection{The Clock}
\label{sec:fundamental_resonance}
\label{sec:fundamental_resonance_filter1}
\textbf{Postulate C.} By Closure the clock is built from the substrate: its period lattice is a nondegenerate plane $\Lambda_c=\mathbb Zu+\mathbb Zv\subset\Lambda$. Its states are the plane's discriminant group $\Lambda_c^*/\Lambda_c$, the finite group determined by the plane and its form alone, and a tick advances through all of them. The clock geometry is determined by its data: all clock tori of the same discriminant are isomorphic.

Because $E_8$ is even, the half-norm $\tfrac12Q(xu+yv)=ax^2+bxy+cy^2$ is an integral binary form of discriminant $D_c=B(u,v)^2-Q(u)Q(v)$ (Lean: \texttt{plane\_form}). The shape $\tau=v/u$ satisfies $a\tau^2+b\tau+c=0$, so the clock torus $\mathbb C/\Lambda_c$ has complex multiplication; this is forced by evenness, not assumed. The discriminant group has exactly $|\det G_c|=|D_c|$ elements for any nondegenerate Gram matrix $G_c$ (Lean: \texttt{card\_discriminant\_group}), so the period is
\begin{equation}
\Delta=|D_c|.
\end{equation}
Tori of discriminant $D_c$ correspond to $SL_2(\mathbb Z)$-classes of forms\cite{silverman_advanced_1994}, so a unique clock geometry means class number one. Over all imaginary quadratic orders this occurs exactly for\cite{stark_complete_1967}
\begin{equation}
\label{eq:heegner_modulus_set}
|D_c|\in\{3,4,7,8,11,12,16,19,27,28,43,67,163\}.
\end{equation}

\label{sec:fundamental_resonance_filter2}
\label{sec:relational_time_constraint}
\textbf{Postulate I.} Each cycle serializes the $N$ channels as active slots, each followed by an idle run. Runs are uniformly bounded, and the schedule is reconstructable from local observations that record only whether an idle run follows a transition, with no external clock or run-length counter.

Gap-only observations reconstruct schedules with runs of length at most $k$ if and only if $k\le1$ (Lean: \texttt{obs\_injOn\_iff}); longer runs alias (\texttt{aliasing}). Hence
\begin{equation}
N\le\Delta\le2N,\qquad 32\le\Delta\le64
\end{equation}
(\texttt{period\_window}). The bound on run length is derived, not imposed: on a single cycle, the periods of reconstructable schedules are exactly those of slot schedules with no two adjacent idle slots (\texttt{reconstructible\_iff\_noAdjacentIdle}), and the scheduling theorems stated with that slot condition as a hypothesis follow with it discharged (\texttt{schedule\_selects\_43\_of\_reconstructible}). The premise carrying the bound is the observation model of Postulate~I: an observation that also recorded run lengths would admit longer runs and remove the upper limit. By Gauss reduction, the only discriminant in this window with a unique positive form is $-43$, proved directly without the classification (\texttt{clock\_unique\_iff}); the vendored class-number-one classification for number fields gives the same window independently. Therefore
\begin{equation}
\Delta=43.
\end{equation}
$E_8$ contains such a plane: $Q(u)=2$, $B(u,v)=1$, $Q(v)=22$, half-norm form $x^2+xy+11y^2$, discriminant group $\mathbb Z/43$ visited in a single cycle (\texttt{e8\_clock\_plane}, \texttt{clock\_glue}). A schedule with eleven idle slots realizes the bound (\texttt{e8\_clock}). Combining the sections so far, the postulates force $E_8$, two $D_4$ sectors, and $\Delta=43$ (\texttt{Vacuum.vacuum\_selects}).

\subsection{Derived Quantities}
\label{subsec:manifold_resolution}
\label{sec:resolution_limits}
\paragraph{Addresses per cycle.} Each of the $D$ coordinates is addressed at each of the $\Delta$ phases:
\begin{equation}
R_M=D\Delta=172.
\end{equation}

\paragraph{Description length.}
\label{par:lembed}
Over one transition a defect is specified by its channel content ($\nu$ coordinates), its configuration ($\sigma$), its boundary closures ($\chi$), and its current and next address ($D$ each, Causality). A description that keeps only some of a state's independent coordinates distinguishes all states exactly when it keeps every one (Lean: \texttt{description\_minimal}); Conservation requires all of them and Economy forbids more. The description therefore has
\begin{equation}
L_{\rm embed}=\nu+\sigma+\chi+2D=31
\end{equation}
coordinates. The substrate gives the same number intrinsically: $E_8$ has $240$ roots and rank $8$, and $(240+8)/8=31$ is its count of nearest-neighbour moves plus coordinate directions per coordinate, the dimension of $\mathfrak e_8$ per rank (\texttt{description\_length}).

All integer inputs are now fixed: $D=4$, $\Delta=43$, $\nu=16$, $N=32$, $\sigma=5$, $\chi=2$, $R_M=172$, $L_{\rm embed}=31$.
